% AY2018 制御工学 第1問〔I〕（転記）
% 出典: 03_制御工学/original/03_制御工学/2018/2018se.pdf
% 要約: 質量 m に並列のバネ k・ダンパ d を介して入力変位 x1(t) を与える系。
%       x2(t) は質量の位置。運動方程式・伝達関数、ステップ応答 x2 と定常値・最大値。

\section*{第1問〔I〕}

図に示す系に関する次の問い (1)〜(4) に答えよ. なお, 質量を $m$, 粘性係数を $d$,
バネ定数を $k$ とし, $x_1(t)$ は系への入力となる変位, $x_2(t)$ は質量の位置を表す.
また, 初期状態において系は静止しており, 粘性に関しては, ダッシュポットの動作速度に
比例した粘性抵抗が生じるものとする.

\begin{center}
\begin{tikzpicture}[>=Latex]
  % 質量 m（左, 変位 x2）
  \draw[thick,fill=gray!12] (0,0.3) rectangle (0.9,1.7);
  \node at (0.45,1.0) {$m$};
  % x2(t)
  \draw[->,thick] (0.45,1.95) -- (1.1,1.95);
  \node at (0.55,2.25) {$x_2(t)$};
  % バネ k（上）
  \draw[thick,decorate,decoration={zigzag,segment length=7,amplitude=4}] (0.9,1.3) -- (2.4,1.3);
  \node at (1.65,1.7) {$k$};
  % ダンパ d（下）
  \draw[thick] (0.9,0.7) -- (1.35,0.7);
  \draw[thick] (1.35,0.5) rectangle (1.8,0.9);
  \draw[thick] (1.6,0.7) -- (2.4,0.7);
  \node at (1.6,0.2) {$d$};
  % 右側連結棒（入力端）
  \draw[thick] (2.4,1.3) -- (2.4,0.7);
  \draw[thick] (2.4,1.0) -- (2.9,1.0);
  \draw[thick,fill=white] (2.9,1.0) circle (0.06);
  % x1(t)
  \draw[->,thick] (2.9,1.0) -- (3.6,1.0);
  \node at (3.2,1.3) {$x_1(t)$};
\end{tikzpicture}
\end{center}

\begin{enumerate}[label=(\arabic*)]
  \item 系に関する運動方程式を求めたうえで, 入力を $x_1(t)$, 出力を $x_2(t)$ として
        伝達関数を求めよ.

  \item 各パラメータを $m=1$, $d=3$, $k=2$ とし, 入力として $x_1(t)$ にステップ入力を
        印加したときの応答 $x_2(t)$ とその定常値を求めよ.

  \item 問い (2) で求めた応答 $x_2(t)$ の最大値ならびにそのときの $t$ を求めよ.

  \item 各パラメータを $m=1$, $d=1$, $k=0$ とし, 入力として $x_1(t)$ にステップ入力を
        印加したときの応答 $x_2(t)$ とその定常値を求めよ.
\end{enumerate}
